\documentclass[11pt]{article}
\usepackage[a4paper,margin=25mm]{geometry}
\usepackage{amsmath,amssymb,amsthm,microtype}
\usepackage[hidelinks]{hyperref}
\newtheorem{theorem}{Theorem}\newtheorem{lemma}[theorem]{Lemma}
\newcommand{\E}{\mathbb E}\newcommand{\Pp}{\mathbb P}
\title{A direct larger-noise comparison from crossing cancellation}
\author{}\date{28 October 2026}
\begin{document}\maketitle

The crossing identity gives a nonnegative loss, but its sign alone does not
pay for the finite binomial and normal corrections. Here is a quantitative
version. A characteristic-function estimate measures the loss for one fixed
count; a coupling transfers it to the selected count. The inward displacement
makes this transfer inexpensive. Only a fixed window of counts then needs a
lower binomial-mass estimate. The rest is bounded by a summable fourth moment.

\begin{theorem}\label{thm:main}
Let $n\ge6000$ be an integer, $p\in[2/5,53/100]$, $q=1-p$, and
$Y=B+U$, with $B$ Bernoulli of parameter $p$ and $\E(U\mid B)=0$.
Assume
\[
 -\tfrac35\le U_0\le a_0,\quad -a_1\le U_1\le\tfrac35,
 \quad a_0,a_1\ge0,\quad d=a_0+a_1\le\min\{9/125,3/n\}.
\]
Put $v=pq$, $s=\E U^2$, $\lambda=ns$, $X=(Y-p)/\sqrt{v+s}$,
$\rho=\E|X|^3$ and $\rho_B=(1-2v)/\sqrt v$. Suppose
$3/40\le\lambda\le9/25$ and $|H|\le46/25$. Write
\[
 t=np+H\sqrt{v+s},\quad k=\lfloor t\rfloor,\quad u=t-k,
 \quad\delta=k-np.
\]
For every $C\ge0$, with $F_B$ the $\mathrm{Bin}(n,p)$ CDF,
\begin{align*}
 \Delta_C={}&\sqrt n\{\Pp(X_1+\cdots+X_n\le H)-\Phi(H/\sqrt n)\}-C\rho\\
 &-\sqrt n\{F_B(k)-\Phi(\delta/\sqrt{nv})\}+C\rho_B
                         <-\lambda/200.
\end{align*}
\end{theorem}

This is a local comparison for all laws in the stated strip. It is not an
unconditional Berry--Esseen cutoff theorem. Its application to a putative
maximizing witness retains the actual-gap acquisition premises. The finite
base and the intermediate indices remain separate obligations.

\section{The exact crossing formula}
Use the independent conditional pairs and increasing mesh from
\texttt{simplification-v13/2026-10-28/crossing.tex}:
\[
 T_j=j+W_j=j+\sum_{i\le j}U_{1i}+\sum_{i>j}U_{0i},
 \qquad T_K\le t<T_{K+1}.
\]
Here $T_0\le t<T_n$: indeed $v\ge6/25$, $\sqrt{v+s}<.501$, and
$|t-np|<.94$. These estimates also give $|\delta|<2$.
Define
\[
 h(x)=\E(x-U_0)_+-\E(x-1-U_1)_+-\Pp(1+U_1\le x),
 \quad R_i=t-(i-1)-\sum_{l<i}U_{1l}-\sum_{l>i}U_{0l}.
\]
The preceding finite telescoping proof gives, with $g_i=\E h(R_i)$,
\begin{equation}\label{eq:age}
 g:=\sum_{i=1}^n g_i=t-\E K
       \ge\E\min\{t-T_K,41/125\},\qquad g_i\ge0.
\end{equation}
In fact $0\le h\le1$, and $h$ vanishes outside $[-3/5,8/5]$.
Below $1-a_1$, it is the increasing function $\E(x-U_0)_+$, whose
limit there is $1-a_1\le1$; above that point it equals
$\Pp(1+U_1>x)-\E(1+U_1-x)_+\le1$.

With $\pi_j=\Pp(B_n=j)$, put
\[
 B_j(t)=\E(t-T_j)_+-(t-j)_+,\qquad
 \mathcal C(t)=\sum_{j=1}^{n-1}(\pi_{j+1}-\pi_j)B_j(t).
\]
The exact weighted formula is
\begin{equation}\label{eq:master}
 \Delta_C=-\sqrt n\sum_i\pi_i g_i+\sqrt n\mathcal C(t)
                    +\mathcal P+C(\rho_B-\rho),
\end{equation}
where
\[
 \mathcal P=\sqrt n\left[u\pi_{k+1}
  -\Phi\!\left(\frac{\delta+u}{\sqrt{n(v+s)}}\right)
  +\Phi\!\left(\frac\delta{\sqrt{nv}}\right)\right].
\]
All lower crossing events are closed and upper crossing events strict.
The identities apply to atoms without a continuity assumption.

\section{Transport from the selected count to a fixed count}
\begin{lemma}\label{lem:transport}
On this one common probability space,
\[
                         \E|W_K-W_k|<13d.
\]
\end{lemma}
\begin{proof}
Let $D_i=U_{1i}-U_{0i}$, and
$V_i=\sum_{l<i}U_{1l}+\sum_{l>i}U_{0l}$.
Centering and the inward caps give $\E|D_i|\le2d$; $V_i$ is
independent of this pair and has variance at most
$n\max(v_0,v_1)\le(5/2)\lambda\le9/10$, where $v_j=\E U_j^2$.

Expand $W_K-W_k$ into $D_i\mathbf1_{\{K\ge i\}}$ for $i>k$
and $-D_i\mathbf1_{\{K<i\}}$ for $i\le k$ and use the triangle
inequality. For $i=k+r$, $r\ge1$, the event $K\ge i$ is
$T_i\le t$, which implies $V_i\le u-r+a_1$.
For $r=1$ use probability at most one; for $r\ge2$ use
Chebyshev with $r-1-d>0$.
For $i=k-r$, $r\ge0$, the event $K<i$ is $T_i>t$, which implies
$V_i>u+r-3/5$. Use probability at most one for $r=0,1$ and
Chebyshev with $r-3/5>0$ thereafter. The endpoint replacements make
these events depend only on $V_i$, so their probabilities multiply
$\E|D_i|$; no stopped-index independence is asserted.

For a decreasing positive function $f$,
$\sum_{r=m}^{\infty}f(r)\le f(m)+\int_m^\infty f(x)\,dx$.
Consequently
\[
 \E|W_K-W_k|\le2d\left[3+\frac9{10}
 \left\{\frac{60}{49}+\left(\frac{125}{116}\right)^2
                        +\frac{125}{116}\right\}\right]<13d.
\]
This bound includes both full tails and is uniform in $n$ and the law.
\end{proof}

\section{A quantitative unweighted loss}
\begin{lemma}\label{lem:cf}
One has $g\ge\lambda/5-13d$.
\end{lemma}
\begin{proof}
For $|x|\le84/125$, the elementary bound
\begin{equation}\label{eq:cos}
                         1-\cos(2\pi x)\ge3x^2
\end{equation}
holds. Here are rational checks avoiding any numerical trigonometry.
On $0\le x\le2/3$, concavity gives
$\sin(\pi x)\ge(3\sqrt3/4)x>(51/40)x$.
On $2/3\le x\le84/125$, the derivative bound $|\sin'|\le1$ gives
$\sin(\pi x)\ge17/20-44/2625$.
Squaring these positive lower bounds proves \eqref{eq:cos}, using
$\sqrt3>17/10$ and $\pi<22/7$.

If $Z$ is either conditional noise and $Z'$ an independent copy, then
$|Z-Z'|\le3/5+d\le84/125$. Thus
\[
 1-|\E e^{2\pi i Z}|^2
       =\E\{1-\cos(2\pi(Z-Z'))\}\ge6\operatorname{Var}Z.
\]
For $x\ge0$ with $2x\le1$, $(1-2x)(1+x)^2\le1$.
It follows that $|\E e^{2\pi i Z}|\le(1+3\operatorname{Var}Z)^{-1}$.
Independence and multiplication give
\[
 |\E e^{2\pi i W_k}|\le\frac1{1+3t_k},\qquad
 t_k=\operatorname{Var}W_k\ge(1-5/n)\lambda\ge\frac{37}{42}\lambda.
\]
The variance estimate follows directly by writing
$t_k=(nq-\delta)v_0+(np+\delta)v_1$ and using $|\delta|<2$,
$p,q\ge2/5$.

For every $r\ge0$, $\sin^2(\pi r)\le\min\{\pi r,1\}$, hence
$\min\{r,41/125\}\ge(3/10)\sin^2(\pi r)$.
Apply \eqref{eq:age} with $r=t-T_K$. Since $K$ is an integer,
\begin{align*}
 g&\ge\frac3{20}\{1-|\E e^{2\pi i W_K}|\}\\
  &\ge\frac3{20}\left(1-\frac1{1+3t_k}\right)
                -\frac{3\pi}{10}\E|W_K-W_k|\\
  &\ge\frac\lambda5-13d.
\end{align*}
In the last step use $3\pi/10<1$, Lemma~\ref{lem:transport}, and
$t_k\ge37\lambda/42$, $\lambda\le9/25$.
The selected-mean bias is already included in \eqref{eq:age}; it has not
been set to zero or paid for a second time.
\end{proof}

\section{Restoring the binomial weights}
For $V_i$ above, independence and centering imply
\[
 \E V_i^4\le(3/5)^2\operatorname{Var}V_i
                +3(\operatorname{Var}V_i)^2
             \le\frac{153}{20}\lambda.
\]
If $i=k+j$, $j\ge11$, then $h(R_i)\ne0$ requires
$V_i\le u+1-j+3/5<13/5-j$.
If $i=k-j$, $j\ge11$, it requires $V_i\ge u+j-3/5\ge j-3/5$.
Since $0\le h\le1$, Markov's inequality and the decreasing-tail integral give
\begin{align*}
 \sum_{|i-k|>10}g_i
 &\le\frac{153\lambda}{20}
       \sum_{j\ge11}\{(j-13/5)^{-4}+(j-3/5)^{-4}\}\\
 &\le\frac{153\lambda}{20}\left[
 (5/42)^4+\frac13(5/42)^3+(5/52)^4+\frac13(5/52)^3\right]
 <\frac\lambda{100}.
\end{align*}

We record explicitly the finite factorial estimate used below:
\[
 \log(m!)=(m+\tfrac12)\log m-m+\tfrac12\log(2\pi)+r_m,
                  \qquad0<r_m<1/(12m),\quad m\ge1.
\]
Its elementary integral-remainder proof is retained in
\texttt{manuscript/arxiv-v1/appendices/count-comparison.tex}, the paragraph
beginning ``For completeness, the finite Stirling formula used here is''.
We use neither a limiting local CLT nor an unspecified asymptotic remainder.

If $|i-k|\le10$, then $|i-np|<12$. For $z=i/n$, every point between
$z$ and $p$ has variance product at least $23/100$. The relative entropy
$D(z\Vert p)$ therefore satisfies
$nD(z\Vert p)\le(i-np)^2/(2n(23/100))$.
The factorial estimate gives
\[
 \sqrt n\pi_i\ge\frac{\phi(0)}{\sqrt{z(1-z)}}
 \exp\left\{-\frac{144/2+1/12}{6000(23/100)}\right\}
 >\frac{199}{250}\left(1-\frac{53}{1000}\right)>\frac34.
\]
Here $z\in[.398,.532]$ and $\phi(0)>199/500$; all bounds are rational
apart from the displayed standard constants, whose enclosures are checked.
Combining this with Lemma~\ref{lem:cf} and $d\le3/n$ gives
\begin{equation}\label{eq:weightedloss}
 \sqrt n\sum_i\pi_i g_i
    >\frac34\left(\frac{19\lambda}{100}-\frac{39}{n}\right)
     =\frac{57\lambda}{400}-\frac{117}{4n}.
\end{equation}

\section{The finite corrections}
The signed-curvature lemma in
\texttt{simplification-v13/2026-10-28/counts-tools.tex} gives
\[
 \sqrt n\mathcal C(t)\le\frac\lambda{4v\sqrt n}
         +\frac{c\sqrt\lambda}{2v\sqrt{n+1}},\qquad
 c=(np-q-t)_+<.46.
\]
Since $v\ge6/25$, this is at most
\begin{equation}\label{eq:curvature}
          \frac{25\lambda}{24\sqrt n}
               +\frac{23\sqrt\lambda}{24\sqrt{n+1}}.
\end{equation}

We next show $\mathcal P<5/n$. Put $\xi=t-np$, so $|\xi|<.94$
and $\delta=\xi-u$. At $z=(k+1)/n$, the segment to $p$ lies in
$[.399,.531]$. On this segment
\[
 |(\log(z(1-z)))'|<1,\qquad |D'''(z\Vert p)|<4.
\]
The factorial estimate and Taylor's theorem consequently give
\begin{equation}\label{eq:massupper}
 \sqrt n\pi_{k+1}\le\frac{\phi(0)}{\sqrt v}
 \exp\left\{-\frac{(\delta+1)^2}{2nv}+e\right\},
 \quad ne\le1+\frac1{12}+\frac{16}{3n}<\frac{11}{10}.
\end{equation}
For $0\le x\le u$, divide this by
$f(x)=\phi((\delta+x)/\sqrt{n(v+s)})/\sqrt{v+s}$.
Its logarithmic upper bound is
\[
 \frac1n\left[\frac{37}{20}
       +\frac{(1-x)(2u+2(47/50)-x-1)}{2v}\right]=:\frac{L(x,u)}n.
\]
The first constant is $11/10+\lambda/(2v)\le11/10+3/4$.
The bracket is nonnegative and less than eight. Also $f(x)<41/50$.
Using $e^z-1\le z/(1-z)$ for $0\le z<1$, integration gives
\begin{align*}
 \sqrt n\pi_{k+1}u-\int_0^u f(x)\,dx
 &\le\frac{41}{50n}\frac{750}{749}
       \left\{\frac{37}{20}+\frac{191}{72}\right\}<\frac{15}{4n}.
\end{align*}
For the integral of the quadratic part, its numerator is
$u[(2-u)(47/50)+2u-(2/3)u^2-1]$, which increases on $[0,1]$
and is at most $47/50+1/3=191/150$.
Finally changing the normal variance at the anchor contributes adversely
only for $\delta>0$, and then contributes less than
\[
 \frac{\phi(0)(47/50)(9/25)}{2(6/25)^{3/2}n}<\frac1n.
\]
This proves $\mathcal P<5/n$ with its full $1/n$ dependence.

The marginal cubic correction is favorable. Let $v_j=\E U_j^2$ and
$\mu_j=\E U_j^3$. Since the two supports have their indicated signs
relative to $p$, the raw absolute third moment is
\[
 v(1-2v)+3v(v_0+v_1)-q\mu_0+p\mu_1.
\]
The one-sided bounds $\mu_0\le a_0v_0$ and
$\mu_1\ge-a_1v_1$ give an increase at least
$(3\min(p,q)-d)s\ge(141/125)s$.
The inequality $(1+x)^{3/2}\le1+3x/2+3x^2/8$ for $x\ge0$
then yields
\[
 (v+s)^{3/2}(\rho-\rho_B)
       \ge s\{87/250-(13/16)s\}>0.
\]
Thus $C(\rho_B-\rho)\le0$ for every $C\ge0$.

Substitute these bounds into \eqref{eq:master}:
\[
 \Delta_C< -\frac{57\lambda}{400}+\frac{137}{4n}
              +\frac{25\lambda}{24\sqrt n}
              +\frac{23\sqrt\lambda}{24\sqrt{n+1}}
           <-\frac\lambda{200}.
\]
For the final inequality divide by $\lambda$, use $n\ge6000$,
$\lambda\ge3/40$, $\sqrt n,\sqrt{n+1}>77$ and
$\sqrt\lambda>27/100$. The resulting strict rational inequality is
included in the companion check. This proves Theorem~\ref{thm:main}.

\section{Why the acquired displacement is small enough}
For an actual central-gap witness, use the notation and the expansion before
AM--GM in \texttt{simplification-v13/2026-10-28/geometry-noise.tex}.
Here $D$ denotes that geometric cost, not the conditional variable above.
On the right strip every noise coefficient in that lower budget is positive;
$Q(p)\ge0$ and the reciprocal-cubic terms are nonnegative. Hence
\[
 nD\ge\frac{6nv\,d(1-d)}{(v+s)^{3/2}}.
\]
If in addition the retained coarse budget gives $nD\le129/4$, then
$n\ge6000$, $\lambda\le9/25$ and $d\le9/125$ imply
\[
 nd\le\frac{(129/4)(v+s)^{3/2}}{6v(1-d)}
       <\frac{(129/4)(251/500)}{6(116/125)}<3.
\]
The actual prior bound $32.20665378$ is smaller than $129/4$.
This is the explicit acquisition join for the new theorem; it does not
replace acquisition, reflection, attainment or the finite base.

For the sharp coefficient, the retained v12 comparison covers smaller
positive noise. Thus within the acquired right strip the new theorem
removes the entire remaining noise range when $n\ge6000$. The v13 geometric
index restriction already excludes acquired $p\le.4$ at such indices.
The range $42\le n<6000$ and the independent finite base remain open.
No unconditional numerical cutoff for the conjecture is asserted.
\end{document}

% Deliberate source-pin mutation.
